\section{Input Data for Surface Interaction Models} \label{sec2.6}

\textbf{General remarks} 

\medskip 

An outline of surface interaction models in general terms was given in
\cref{sec1.3}.
As for the implementation of such models in EIRENE, there are two parts to the
surface interaction data block.
The first part contains data which are general to the EIRENE reflection model
(``Block for General Reflection Data"), at all surfaces.
The second part may be different for each individual surface element (label:
$MSURF$) and thus must be specified for each reflecting ``non-default" surface
of the ``standard mesh" and for each reflecting ``additional surface".
It is referred to as ``Block for Local Reflection Data".
If this block is missing and if the surface $MSURF$ is neither transparent
($ILIIN<0$) nor purely absorbing ($ILIIN=2$) nor a ``mirror surface"
($ILIIN=3$) nor a ``periodicity surface", ($ILIIN\ge 4$) then the default
reflection model (i.e.\ $Fe-Target$, recycling coefficient = one, etc.) is
activated for this surface.

The surface interaction model comprises the following information:

\begin{enumerate}
  \item $p_f$ (RPROBF): probability for reflection of a ``fast" particle

        (= 0.
        for incident molecules),
  \item $p_t$ (RPROBT): probability
        for re-emission as ``thermal particle"

        (with surface temperature, see the flags EWALL)
  \item $p_a$ (RPROBA):
        probability for absorption 

        (surface sticking, pumping, etc.)
  \item $p_s$ : probability for (physical) sputtering
  \item $p_c$ : probability
        for (chemical) sputtering
  \item For the reflected ``fast particles" and for the
        sputtered particles the type, species and the distribution of energy, polar and
        azimuthal angle of emitted particles must be specified.
  \item For the re-emitted thermal particles the type, species and only the
        distribution of energy must be given; the angular distribution follows the
        cosine law by default.
\end{enumerate}

All these data may be functions of incident energy $E_{in}$, of incident polar
angle $\theta_{in}$ against the surface normal and of projectile and target
species.
The three probabilities $p_f$, $p_t$ and $p_a$ must add up to 1 (one); $p_s$
and $p_c$ are sputtering yields per incident particle and, thus, may
occasionally be larger than 1.
They are called a probability here only by abuse of language.

All this information is contained in each of three so called ``Reflection
Models", namely the ``Database Reflection model" (see references
\cite{kn:Eckstein}, \cite{kn:Bateman}), the ``Modified Behrisch Matrix model"
(references \cite{kn:Behrisch}, \cite{kn:Nicolai}) or the ``User Supplied
Reflection model" (see \cref{sec3.3}), and in several so called ``Sputter
Models" (\cite{kn:NucFus}, \cite{kn:Sputer95}, \cite{kn:Roth96}.

These models may, to some extend, be modified by the input flags described
here.
For example, a surface may be purely absorbing (for all test particles incident
onto it, $ILIIN = 2$ option), acting like a mirror for the neutral test
particles, by the $ILIIN = 3$ option (see section 2.3A and 2.3B), or enforce
periodicity, e.g., because of symmetry conditions ($ILIIN \ge 4$ option).
In these two latter cases all test particles of all species are reflected with
the species unchanged, with probability one and elastically:
$E_{out} = E_{in}$.
In case of the ``mirror reflection option" the normal component $\vvv_n$ of the
incident velocity vector $\vvv$ is reversed: $\vvv_n \rightarrow -\vvv_n$.
In case of a ``periodicity surface" both the position and the speed unit vector
are altered according to the specific periodicity of the configuration.

If, instead, a surface reflection model is chosen ($ILIIN = 1$), then some of
the input flags described here (RECYCF, RECYCT, RECYCS, RECYCC) as well as the
flags (RINTEG, EINTEG, AINTEG) may be used to modify conveniently the otherwise
rather unhandy reflection data and formulas.

\medskip

\textbf{The Input Block for General Reflection Data} 

\begin{lstlisting}
*** 6A. Data for General Reflection Models
      NLTRIM
C     ``Path-Card'' for TRIM surface reflection files
C     (machine specific, may be omitted)
\end{lstlisting} 

If a ``Path Card'' is included, then read an arbitrary number of
``Target-Projectile Specification Cards'' of the format 

\begin{lstlisting}
      A_on_B
\end{lstlisting} 

Up to NHD6 (see PARAMETER statements, \cref{sec3.1}) data files for different
Target-Projectile combinations can be read in EIRENE versions older than 2001.
No such limitation exists in more recent versions (due to dynamic allocation of
storage).

If no ``Path Card'' has been specified, or after the ``Target-Projectile
Specification Cards", continue reading general species sampling distributions,
used for source particle species sampling as well as for sampling post-surface
event species as part of the surface reflection models 

\begin{lstlisting}
      (DATD(I) I=1,NATM)
      (DMLD(I) I=1,NMOL)
      (DIOD(I) I=1,NION)
      (DPLD(I) I=1,NPLS)
C (this card is to be read only in case NPHOT > 0)
      (DPHT(I) I=1,NPHOT) 
      ERMIN  ERCUT  RPROB0  RINTEG EINTEG AINTEG
\end{lstlisting} 

Note that the species index distribution cards must always be read, even if
NATM, NMOL, NION or NPLS are zero.
In this latter case at least one (irrelevant) parameter must be read.
E.g.\ if NION=0, i.e. no test ion species in particular run, then still one
card with the irrelevant parameter DIOD(1) must be read.
An exception from this rule is the photon species distribution DPHT: this card
is only read if there are ``photonic test particles" in the run, i.e.\ if NPHOT
$>$ 0.
This exception was necessary for backward compatibility of input file format
after introducing photons as new type of species, around 2001.

Next, optionally, an arbitrary number of surface models, labelled by the
character string \lstinline|modname|, may be read.

\begin{lstlisting}
     SURFMOD_modname
     ILREF  ILSPT ISRS  ISRC
     ZNML EWALL EWBIN TRANSP(1,N) TRANSP(2,N) FSHEAT
     RECYCF  RECYCT  RECPRM  EXPPL  EXPEL  EXPIL
C (following line may be omitted, then: default sputter 
C  model, see below)
     RECYCS  RECYCC  SPTPRM  ESPUTS ESPUTC
\end{lstlisting}

next: within each SURFMOD-sub-block there may be an
arbitrary number of lines

\begin{lstlisting}
      VARNAME  SPZNAME   VALUE
C     (format: CHARACTER*8,X,CHARACTER*8,*)
\end{lstlisting}

\medskip

\textbf{Meaning of the Input Variables} 

\medskip 

\begin{description}
  \item[NLTRIM] TRIM database is used, if ``Database
        Reflection Model" is specified in at least one block for local reflection data.
        Data are read from data-set FT21 (no ``Path Card" specified, old option), or
        from the domain specified by the ``Path Card" (new option, see next card).
        If the old option is used, then the complete TRIM file is read, containing the
        first 12 TRIM target-projectile combination data-sets listed in \cref{sec1.3},
        i.e., the files H\_on\_Fe to T\_on\_W.

        If the parameter NHD6 $<$ 12, (\cref{sec3.1}) then only the first NHD6 files
        are read from FT21.
  \item[``Path Card"] This card is machine specific.
        If EIRENE finds a card containing the string 'PATH' or 'path', it assumes that
        this card specifies the path to the domain containing the TRIM surface
        reflection data files.
  \item[A\_on\_B] Name of a particular TRIM data file in the domain specified by
        the path card.
        E.g., H\_on\_Fe would include the data file for hydrogen onto iron into the
        EIRENE run.
        Up to NHD6 such ``Target-Projectile Specification Cards" may be included.
        For a complete list of such files currently available see again \cref{sec1.3}.
\end{description}
Note: if during a Monte Carlo simulation a projectile A hits a target B, for
which no TRIM data file has been specified, but still NLTRIM=TRUE, then EIRENE
searches the ``closest" of all its TRIM data files (with respect to a reduced
mass argument) and uses this target-projectile combination together with a
reduced mass scaling of incident particle energy.
Hence, a TRIM data set is chosen such that the reduced mass scaling factor is
as close to one as possible amongst the files available.
\begin{description}
  \item[DATD] distribution for sampling the species index of reflected or
        otherwise emitted atoms.
        The NATMI relative frequencies 

        DATD(IATM) IATM = 1,NATMI 

        are used to produce the corresponding cumulative distribution DATM in order to
        facilitate sampling (inversion method).
        Normalization of DATD such that 

        $\sum_{IATM} DATD(IATM) =1$ 

        is carried out internally.
  \item[DMLD] as for DATD, but for molecules.
        Cumulative distribution is DMOL.
  \item[DIOD] as for DATD, but for test ions.
        Cumulative distribution is DION.
  \item[DPLD] as for DATD, but for bulk ions.
        Cumulative distribution is DPLS.
\end{description}
The activation of these distributions at surface events during the particle
history generation process is controlled by the surface species index flags
ISRF\$, ISRT\$ read in the species sub-blocks of block 4 and 5
(\cref{sec2.4,sec2.5}).
By default the following conventions are used (and can be overruled by calls to
subroutines REFUSR, SPTUSR only, see \cref{sec3.3}):
\begin{description}
  \item[ISRF\$] fast particle reflection species flag:

        \begin{description}
          \item[$> 0$] If ISRF\$ $\le$ NATMI, then ISRF\$=IATM, the species labelling
                index for the reflected fast atom.

                If ISRF\$ $>$  NATMI not in use, warning and error exit.
          \item[$= 0$] no fast particle reflection for this species: $p_f = 0 $
          \item[$<
                  0$] not in use, warning and error exit.

        \end{description}
  \item[ISRT\$] thermal particle re-emission species flag:

        \textbf{incident atoms, test ions and bulk ions:}
        \begin{description}
          \item[$>
                  0$] If ISRT\$ $\le$ NATMI, then ISRT\$ = IATM, the species labelling index for
                the re-emitted thermal atom.

                If ISRT\$ $>$  NATMI, then IATM is sampled from the distribution DATM(IATM).
          \item[$= 0$] no thermal particle re-emission for this species: $p_t = 0 $
          \item[$< 0$] molecule is re-emitted, and -ISRT\$ is used to identify the
                species labelling index IMOL for the re-emitted molecule, exactly as described
                above for re-emitted atoms.
                The relevant sampling distribution for species index IMOL in case -ISRT\$ $>$
                NMOLI is DMOL(IMOL).
                This option can e.g.\ be used to specify the distribution of vibrationally
                excited molecules $H_2(\nu)$ emitted from a surface after bombardment by
                incident $H$ atoms, the so called Eley Rideal mechanism:
                Surface(with implanted $H) + H \rightarrow H_2(\nu)$
        \end{description}
        \textbf{incident molecules:}
        \begin{description}
          \item[$> 0$] If ISRT\$ $\le$
                NMOLI, then ISRT\$ = IMOL, the species labelling index for the re-emitted
                thermal molecule.

                If ISRT\$ $>$  NMOLI, then IMOL is sampled from the distribution DMOL(IMOL).
          \item[$= 0$] no thermal particle re-emission for this species: $p_t = 0 $
          \item[$< 0$] not in use, warning and error exit:
                ``\lstinline|species index out of range|'' 

                In code versions older than year 2000 this ISRT$<$0 options was also used to
                activate sampling the new species index IMOL from distribution DMOL.
                In order to avoid un-intentional species sampling at surfaces this option was
                disabled from year 2000 on.
        \end{description}
\end{description}
These latter two surface-species index flags are considered to be particle
properties, hence they are read in the particle specification blocks 4 and 5
(\cref{sec2.4,sec2.5}).
The flags in the next card can be used to modify the preprogrammed fast
particle reflection models, to some extend at least.
\begin{description}
  \item[ERMIN] For incident particle energies below ERMIN, the ``fast" particle
        reflection model is switched off.
        Only the ``thermal" particle model is used.
  \item[ERCUT,RPROBF] ~ 

        These variables may be used to modify the default ``Behrisch Matrix" reflection
        coefficients for particles incident on a surface at low energies $E_{in}$.
        The original data \cite{kn:Behrisch} are used only for $E_{in}$ $>$ ERCUT and
        for normal incidence $\vartheta_{in} = 0 $.

        In the range ERMIN $<$ $E_{in}$ $<$ ERCUT the Behrisch Matrix particle
        reflection coefficient $p_f(E_{in}, \vartheta_{in} = 0)$ is replaced by a
        smooth cubic interpolation curve $p_f (E_{in})$ such that $p_f(E_{in},
          \vartheta_{in} = 0)$ = RPROBF .

        The original ``Behrisch Matrix" is recovered by setting ERCUT $\le 0$.
\end{description}
The next three flags modify the particle, energy and momentum reflection
coefficients, and/or the resulting distributions in post-reflection energy and
angle, respectively.
These flags only apply to incident particles of type 1, 3, 4 and 5, i.e., not
for incident molecules because for those RPROBF = 0 by default.
\begin{description}
  \item[RINTEG] ~
        \begin{description}
          \item[$> 0$] Fixed (independent of energy
                and angle of incidence) particle reflection coefficient.

                The fast particle reflection probabilities RPROBF are set to $p_f = MIN(1 -
                  p_a, RINTEG)$, regardless of the reflection model selected by the flag ILREF in
                block 6B.
                $p_a$ is kept as specified, and the thermal component probability $p_t$ is then
                recomputed as $p_t = 1 - p_f - p_a$.

                RINTEG $\ge 1 - p_a$ enforces the fast particle reflection model for all
                un-pumped incident particles, i.e., RPROBF = RINTEG is internally reset to $1 -
                  p_a$.
          \item[$= 0$] Default: fast particle reflection probability as defined by the
                reflection model chosen.
          \item[$< 0$] The fast particle reflection probabilities RPROBF are set to
                $1.0-p_a$.
                Hence: same as RINTEG=1.0
        \end{description}

  \item[EINTEG] ~
        \begin{description}
          \item[$> 0$] Fixed (independent of energy
                and angle of incidence) energy reflection coefficient.

                This choice replaces the energy random sampling procedure in the fast particle
                reflection model by an reflection assumption:
                $E_{out} = E_{in} \cdot EINTEG$.

          \item[$= 0$] Default: no modification of energy distribution in the reflection
                model for fast particle reflection.
          \item[$< 0$] elastic ($E_{out} = E_{in}$) reflection for all particles
                reflected according to the fast particle reflection model.
                Hence: same as EINTEG=1.0
        \end{description}
  \item[AINTEG] ~
        \begin{description}
          \item[$> 0$] 

                Fixed (independent of energy and angle of incidence) momentum reflection
                coefficient (``accommodation coefficient'') $\alpha = 1-AINTEG$

                This choice replaces the angle random sampling procedure in the fast particle
                reflection model by activating features of the so called ``Maxwell's boundary
                model'' (Philosophical Transactions of the Royal Society, 1879), in which a
                fixed fraction AINTEG ($\leq 1.0$) is reflected at a specular angle.
                The remaining ``evaporated'' fraction (1-AINTEG) is reflected with a cosine
                distribution (Lambertian).
                In the original Maxwell boundary model the evaporated fraction is re-emitted
                according to an inward directed Maxwellian flux distribution (taken at wall
                temperature) while the specular fraction is with unchanged energy $E_{in} =
                  E_{out}$.
                This full Maxwell boundary condition with accommodation coefficient
                $\alpha=1-$AINTEG is specified with the additional choices:
                EINTEG=1.0, EWALL =-TW (TW the wall temperature, see further below: EWALL flag.

          \item[$= 0$] Default:
                no modification of angular distributions in reflection model
                for
                fast particle reflection.
          \item[$< 0$] specular reflection for all particles reflected according to the
                fast particle reflection model.
                Hence: same as AINTEG=1.0
        \end{description}
\end{description}

\textbf{Note:} 

Some ``General Reflection Model" data of input block 6A are copied identically
NLIMI + NSTSI times onto appropriately dimensioned arrays (with the same names)
in order to make them dependent upon the surface labelling index as well.

The array index is then the surface number.

Presently this ``localization option" is available for the flags

\begin{lstlisting}
      RINTEG, EINTEG, AINTEG
\end{lstlisting} 

Overwriting general reflection data for some specific surfaces can be done via
a call to the entry RF0USR of the user supplied reflection routine REFUSR.
RF0USR is called in the initialization phase of an EIRENE run (see
\cref{sec3.3}).

By a similar strategy, some of the species independent flags in the next
sub-block 6B can be made species-dependent.
See, for example, RECYCF, RECYCT, RECYCS and RECYCC below.

\medskip

\begin{lstlisting}
*** 6B. Data for Local Reflection and Sputtering Models}
\end{lstlisting}

\medskip

The next 3 (or 4) lines comprise the sub-block for local reflection data.
Such sub-blocks can be included in each ``surface deck" (block 3A, 3B) to
overwrite the default reflection model for each individual surface.

\begin{lstlisting}
      ILREF  ILSPT  ISRS  ISRC
      ZNML EWALL EWBIN TRANSP(1) TRANSP(2) FSHEAT
      RECYCF  RECYCT  RECPRM  EXPPL  EXPEL  EXPIL
C (following line may be omitted then: default sputter model, 
C  see below)
      RECYCS  RECYCC  SPTPRM  ESPUTS ESPUTC
\end{lstlisting}

\medskip

An arbitrary number of such sub-blocks of 3 (or 4) lines each may also be
defined under a certain label \lstinline|SURFMOD_modname| and be included in the
input file in block 6 directly after the ERMIN, \dots\ deck 

In blocks 3a and 3b surfaces may be assigned a particular local reflection
model by a card reading \lstinline|SURFMOD_modname|.
Many surfaces may then be linked to the same surface reflection sub-block.

Example:

\begin{footnotesize}
  \begin{lstlisting}
SURFMOD_BERYL_SPT_300K
     1     2     0     0
 9.04000E+02-2.60000E-02 0.00000E+00 0.00000E+00 0.00000E+00 2.80000E+00
 1.00000E+00 1.00000E+00 1.00000E+00 1.00000E+00 5.00000E-01 1.00000E+00
 1.00000E+00 1.00000E+00 1.00000E+00 0.00000E+00 0.00000E+00
SURFMOD_CARB_SPT_1153K
     1    12     0     0
 1.20600E+03-1.00000E-01 0.00000E+00 0.00000E+00 0.00000E+00 2.80000E+00
 1.00000E+00 1.00000E+00 1.00000E+00 1.00000E+00 5.00000E-01 1.00000E+00
 1.00000E+00 1.00000E+00 1.00000E+00 0.00000E+00 0.00000E+00
TRANSP1  D2      0.5000E 00
TRANSP2  T2      0.2500E 00
RECYCF   D       0.0000E 00
RECYCT   D2      0.9500E 00
SURFMOD_CARB_SPT_812K
     1     2     0     0
 1.20600E+03-7.00000E-02 0.00000E+00 0.00000E+00 0.00000E+00 2.80000E+00
 1.00000E+00 1.00000E+00 1.00000E+00 1.00000E+00 5.00000E-01 1.00000E+00
 1.00000E+00 1.00000E+00 1.00000E+00 0.00000E+00 0.00000E+00
\end{lstlisting}
\end{footnotesize}

\textbf{Meaning of the Input Variables} 

\medskip 

\begin{description}
  \item[ILREF] Flag for choice of local reflection model
        \begin{description}
          \item[$= 1$] TRIM database reflection model is used.
                NLTRIM must be .TRUE.
          \item[$= 2$] ``modified Behrisch Matrix model" is used
          \item[$= 3$] user
                supplied reflection model (see \cref{sec3.3}:
                Subroutine REFUSR)
        \end{description}

        Default: ILREF = 2
  \item[ILSPT] Flag for choice of local sputtering model

        Let ILSPT = MN, with M and N single digit integers each.
        Then N controls the options for physical sputtering, and M controls chemical
        sputtering.
        See subroutine SPUTER.
        \begin{description}
          \item[N = 0] no physical sputtering at this surface
          \item[N = 1] constant
                physical sputtering rate (see parameter RECYCS below)
          \item[N = 2] modified
                Roth--Bogdansky formula for sputter yield, Thompson energy distribution and
                cosine angular distribution for emitted particles (see references
                \cite{kn:Sputer95} and \cite{kn:Roth96}).
          \item[N = 9] (was option N=3 in Eirene-2004 and older) 

                user supplied sputtering model (see \cref{sec3.3}: Entry SPTUSR to subroutine
                REFUSR)
          \item[M = 0] no chemical sputtering at this surface
          \item[M = 1]
                constant chemical sputtering rate (see parameter RECYCC below)
          \item[M = 2]
                ``Roth formula" for chemical sputter yield, thermal distribution for emitted
                particles (see reference \cite{kn:Roth99}), ``weak flux dependence option A6".
          \item[M = 3] ``Roth formula" for chemical sputter yield, thermal distribution
                for emitted particles (see reference \cite{kn:Roth99}), ``strong flux
                dependence option A7".
          \item[M = 4] ``Roth formula" for chemical sputter yield, thermal distribution
                for emitted particles (see reference \cite{kn:Roth99}), ``new flux dependence
                option A8 (2004)".
          \item[M = 5] not in use
          \item[M = 6] ``Haasz--Davis 1998 formula" for chemical
                sputter yield
          \item[M = 7] ``Haasz--Davis 1998 formula" for chemical sputter
                yield, and multiplicative factor for flux dependence (Roth, 2004).
          \item[M = 9] (was option N=3 in Eirene-2004 and older) 

                user supplied sputtering model (see \cref{sec3.3}: Entry SPTUSR to subroutine
                REFUSR)
        \end{description}

        Default: ILSPT=0
\end{description}
The next two surface-species index flags
ISRS\$ and ISRC\$ control the species of sputtered particles.
Hence they are considered surface properties and are read in the surface
specification decks in blocks 3a, 3b and/or 6 (\cref{sec2.3a,sec2.3b,sec2.6}).
\begin{description}
  \item[ISRS\$] sputtered particle species flag (physical sputtering).

        \begin{description}
          \item[$> 0$] both the sputtered particle and the reflected particle (if any)
                will be followed.
                Their contribution to surface particle and energy fluxes is stored in surface
                averaged tallies 1 to 24, i.e., sputtered particles are not explicitly
                distinguished from reflected particles in the particle and energy balances.
                Furthermore the ``sputtered flux surface tallies" 33 to 37
                (\cref{sec.tables-old}) in older versions before 2002, and on tallies 51 to 81
                (\cref{sec.tables-new}) else,  are also updated.
                The species index of the sputtered particle (atom) is IATM=ISRS.
                Hence, on input, one must have 1 $\le$ ISRS $\le$ NATMI, otherwise: error exit.

          \item[$\leq 0$] Same as ISRS $>$ 0, however, the species index of the sputtered
                particle is determined automatically from comparing the charge and mass numbers
                of the available atomic species (input block 4a) with the corresponding surface
                material (nuclear mass and charge) data of the surface element.
                If no suitable atomic test particle is is found, then only sputter tallies are
                scored, but no sputtered particles are then subsequently traced.
          \item[$= 0$] Same as ISRS $\leq$ 0, but in this case a sputtered particle is
                NOT followed, even if its atomic test particle species could be identified.

                Only the reflected particles are followed.
                and only their contribution to surface particle
                and energy fluxes is stored in regular surface averaged tallies
                1 to 24.
                Sputter tallies are still scored.

        \end{description}

        Note: ISRS=ISRS(ISPZ,MSURF), so the above described options for sputtered
        particle species, as well as for either only scoring fluxes or even tracing
        these sputtered particles, can be made dependent on the incident species index
        ISPZ.

  \item[ISRC\$] sputtered particle species flag (chemical sputtering).

        \begin{description}
          \item[$> 0$] both the sputtered particle and the reflected particle (if any)
                will be followed.
                Their contribution to surface particle and energy fluxes is stored in surface
                averaged tallies 1 to 24, i.e., sputtered particles are not explicitly
                distinguished from reflected particles in the balances.
                Furthermore the ``sputtered flux surface tallies" 25 to 28 are updated.
                The species index of the sputtered particle (atom) is IATM=ISRC, if ISRC $\le$
                NATMI, or (molecules) IMOL, if ISRC = NATMI+IMOL and NATMI $<$ ISRC $\le$
                NATMI+NMOLI.
                Hence, on input, ISRC $\le$ NATMI+NMOLI.
          \item[$= 0$] There is no chemical sputtering for the particular surface element
                and incident species (note: ISRC=ISRC(ISPZ,MSURF), i.e., $p_c = 0$ here.

                Only the reflected particles are followed.
                Their contribution to surface particle and energy fluxes is stored in surface
                averaged tallies 1 to 24.
          \item[$< 0$] Same as ISRC $>$ 0, however, the species index of the sputtered
                particle is determined automatically from comparing the charge and mass numbers
                of the atomic species (input block 4a) with the corresponding data of the
                surface element.
                I.e., in case of Carbon surfaces the sputtered particle is a C-atom, if such an
                atom has been specified in input block 4a
        \end{description}

  \item[ZNML] = KLMN (4 digits)
        \begin{description}
          \item[KL] atomic weight of
                wall material.
                %\end{description}
                Note: the nearest integer of the mass number in the TRIM runs is used.
                For example, a copper target is specified in the TRIM files with an atomic
                weight of 63.54, and the corresponding TRIM file is used for surfaces with
                KL=64.
          \item[MN] nuclear charge number of wall material
        \end{description}
        Example:
        Carbon: ZNML=1206.

        Example: Molybdenum: ZNML=9642.

        Example: Copper: ZNML=6429.

        Default: ZNML = 5.626E3 (stands for Fe).
  \item[EWALL] ~
        \begin{description}
          \item[$< 0$] -EWALL = TW is a (surface-)
                temperature (\si{\electronvolt}) in a Maxwellian flux distribution for the
                thermal particle energy.
                The resulting mean energy is $E_{mean} = 2 \cdot TW$.
          \item[$> 0$] +EWALL = Energy of mono-energetic (thermal) particles.

                The relation between surface temperature TW and the mean energy of particles
                then reads EWALL = $E_{mean} = 1.5 \cdot TW$.
          \item[$= 0$] Energy is sampled from a Thompson distribution, using the flag
                EWBIN (see below) as parameter for the surface binding energy
        \end{description}

        Default: EWALL = +0.0388 ($\simeq$ TW = \SI{0.026}{\electronvolt} $\simeq$
        \SI{300}{\kelvin}) 

        Note that the $EWALL > 0 $ option enables EIRENE to include boundary conditions
        in ``one speed transport equation'' approximations, which often are of great
        interest in general linear transport theory.
        The $EWALL < 0 $ option, together with the EINTEG and AINTEG flags described
        above, provides the Maxwell boundary conditions with given wall accommodation
        coefficient (see above).
  \item[EWBIN] see above, EWALL = 0 option.

        Default: EWBIN = 0.0 (irrelevant for ``Default Model'') 

  \item[TRANSP(1)] Semi-transparency for particles incident from the positive
        side.
        Renders a non-transparent surface (ILIIN $>$ 0) semi-transparent.
        The probability of a test flight to pass through the surface is [TRANSP].
        This transparent fraction of the current is not scored on surface tallies,
        i.e.\ for this fraction of surface currents the ILIIN=0 option is used, see
        input block 3B.
        Hence:
        the probability for reflection/re-emission/absorption according to
        ILLIN setting is [1-TRANSP].
        Only this fractions [1-TRANSP] of the surface currents are scored as ``incident
        flux" and, if applicable, ``re-emitted flux" surface tallies.
        In order to also score the net current of the transparent flux fraction
          [TRANSP] on surface tallies, an additional ``very nearby" fully transparent
        scoring-surface can be used.

        Default: 0.0  (i.e., fully reflecting surface).

        Irrelevant for transparent surfaces (with ILIIN $\leq 0$).

  \item[TRANSP(2)] Semi-transparency for particles incident from the negative
        side on a non-trans\-parent surface.
        Otherwise same as for TRANSP(1) option.

        Default: 0.0  (i.e., fully reflecting surface).

        Irrelevant for transparent surfaces (with ILIIN $\leq 0$).

  \item[FSHEAT] surface sheath potential factor.
        The sheath potential is $FSHEAT \cdot T_e$, with $T_e$ the electron temperature
        at the point of incidence.
        This sheath potential is applied if ions (test ions or bulk ions) hit a
        non-transparent surface.

        If FSHEAT $\le$ 0.0, then a sheath potential computed from the local background
        plasma flow conditions is used (function SHEATH), assuming ambipolar flow, a
        Boltzmann distribution for electrons and zero secondary electron emission.
        See \cref{sec1.4}.
        In case of zero (undefined) background plasma flow velocity at the place of
        incidence, a default of FSHEAT = 2.8 is used (corresponding to $T_e = T_i$, and
        a single ion species $D^+$ plasma flowing at ion acoustic speed parallel to the
        $\vec B$-field into the sheath.

        Default: FSHEAT = 0.0 

  \item[RECYCF, RECYCT] Multiplier for reflection probability RPROBF:

        Particles can be re-emitted from surfaces by either the ``fast reflection''
        model or by a ``thermal emission'' model.
        Flag RECYCF controls (scales) the ``fast particle reflection" probability
        $p_f$.

        The probability $p_f = RPROBF(E_{in},\Theta_{in},ispez,wall)$ for the ``fast''
        particle reflection model, as specified by other flags for this surface, is
        modified to
        \begin{equation}
          \label{rprobf}
          RPROBF(E_{in},\Theta_{in},ispez,wall) = AMIN( RECYCF \cdot RPROBF, RECYCT),
        \end{equation}
        where RPROBF was evaluated from the reflection model specified
        by ILREF.
        Note the cut-off at recycling coefficient RECYCT (defined below).

        The total recycling coefficient RECYCT = $p_f + p_t$ is unchanged by flag
        RECYCF.
        Hence, by the use of RECYCF not only the fast particle reflection probability,
        but also the probability for thermal particle emission $p_t$ is altered to
        maintain a total recycling coefficient at this surface of RECYCT.

        Default: RECYCF = 1 for incident atoms, test ions and bulk ions.

        Default: RECYCF = RPROBF =0 for incident molecules.
  \item[RECYCT] Recycling coefficient (must not be negative):

        A flux $RECYCT \cdot Influx$ is re-emitted from a surface, for any $Influx$ of
        particles of any species, where all fluxes are measured as ``atomic fluxes"
        (=fluxes of nuclei).
        RECYCT hence defines the sticking probability $p_a$ [and hence also the pumping
            speed \cref{pumping-speed}] of any surface in EIRENE, for all incident species.

        The fraction $p_a$ = (1 - RECYCT) of incident (atomic) flux will be absorbed at
        the surface.
        The non-sticking, i.e.\ the re-emitted fraction $RECYCT=1-p_a$ is split into a
        ``fast" and a ``thermal'' component.

        A fraction $p_f$ = RPROBF [see \cref{rprobf}] of the incident particles is
        reflected as described by the ``fast particle reflection model".
        However, by relation \cref{rprobf} it is ensured that RPROBF is always less
        than or equal to RECYCT.

        The fraction $p_t$ = RPROBT = (RECYCT - RPROBF) will be re-emitted by the
        ``thermal particle reflection model".
        This flag is to be used to define an effective pumping speed at certain
        surfaces, see \cref{sec2.6.1} below.

        Note: RPROBF = 0 for incident molecules (ITYP=2) by default.

        Default: RECYCT = 1., i.e.\ $p_a = 0$
  \item[RECPRM] free model parameter for
        user supplied recycling models $ILREF=9$.

        Default: RECPRM = 0.

  \item[EXPPL] (only for ILREF = 2 option) 

        incident angular dependence of fast particle reflection coefficient.
        The formula
        \begin{displaymath}
          R( \Phi ) = 1 - (1 - RPROBF)	\cdot
          \cos^{EXPPL} (\Phi)
        \end{displaymath}
        is used, $e1 = EXPPL$, see \cref{eq1.18},
        where:
        \begin{description}
          \item[RPROBF] Reflection probability from ``Behrisch Matrix" model, which is
                valid only for normal incidence.
          \item {$\Phi$ \hfill} Angle of incidence against surface normal
          \item
                {R($\Phi$) \hfill} Reflection probability for particles incident with angle
                $\Phi$ 

                note: $R(\Phi) = 1$ for $\Phi = 90^{\circ}$ and EXPPL $>$  0.
        \end{description}
        Default: EXPPL = 1.
        (recommended from a comparison with the TRIM database)
  \item[EXPEL] (only for
        ILREF = 2 option) 

        as EXPPL, incident angular dependence for the energy reflection coefficient.

        Default: EXPEL = 0.5 (recommended from a comparison with the TRIM database).
        e2 = EXPEL, see \cref{eq1.19,eq1.20}
  \item[EXPIL] (only for ILREF = 2 option) 

        incident angular dependence for angular distribution of re-emitted atom or
        molecule; 

        \begin{description}
          \item[$= 0$] cosine distribution (Lambertian)
          \item[$> 0$]
                mixed cosine-specular model.

                the specular contribution increases according to \cref{eq1.21}
                with angle of incidence $\Phi$ and with
                e3 = EXPIL (recommended: $EXPIL \le 1$).
          \item[$> 100$] specular reflection angles
        \end{description}

        Default: EXPIL = 0.

  \item[RECYCS] The meaning of this flag for the ``physical sputtering" options
        depends upon the value of the first digit N of ILSPT:
        \begin{description}
          \item[N = 0] no physical sputtering, YIELD1 = 0.
                RECYCS is irrelevant.
          \item[N = 1] constant physical sputtering yield, YIELD1 = RECYCS
          \item[N = 2]
                RECYCS is a multiplier for the sputtered particle flux YIELD1.
                YIELD1 is computed from the incident species, energy, angle and surface
                parameters by the sputter model $N=2$.
                Hence:
                the sputtered particle yield YIELD1
                as computed from subroutine SPUTER
                is modified to

                $YIELD1 = RECYCS \cdot YIELD1 $.
          \item[N = 9] RECYCS is a free model parameter, which can be used in the user
                supplied sputter model for any particular surface element.

                Default: RECYCS = 1.
        \end{description}
  \item[RECYCC] The meaning of this flag for the ``chemical sputtering" options
        depends upon the value of the second digit M of ILSPT:
        \begin{description}
          \item[M = 0] no chemical sputtering, YIELD2 = 0.
                RECYCC is irrelevant.
          \item[M = 1] constant chemical sputtering yield, YIELD2 = RECYCC.
          \item[M = 2] RECYCC is a multiplier for the sputtered particle flux YIELD2.
                YIELD2 is computed from the incident species, energy, angle and surface
                parameters by the sputter model $M=2$.
                Hence:
                the sputtered particle flux YIELD2
                as computed from subroutine SPUTER
                is modified to

                $YIELD2 = RECYCC \cdot YIELD2 $.
          \item[M = 9] RECYCC is a free model parameter, which can be used in the user
                supplied sputter model for any particular surface element.

                Default: RECYCC = 1.
        \end{description}
  \item[SPTPRM] free model parameter for user supplied sputtering models N=3,
        M=3.

        Default: SPTPRM = 0.
  \item[ESPUTS] (new: March 2015) parameter (flag) for energy of physically
        sputtered particle.
        Currently not in use.

        Default: ESPUTS = 0.
  \item[ESPUTC] (new: March 2015) parameter (flag) for energy of chemically
        sputtered particle.
        By default: chemically sputtered particles are released from the wall by the
        ``thermal surface emission model", as also used for the recycling/reflection
        thermal emission, i.e.\ with wall temperature (as specified by EWALL parameter,
        see above) and either a cosine (for monoenergetic emission at EWALL = 1.5
        TWALL) angular distribution, or sampling from a stationary Maxwellian flux
        distribution at -EWALL = TWALL.

        If ESPUTC .GT.\ 0, then the the chemically sputtered particles are released
        with a mono-energetic distribution at E0 = ESPUTC, and a cosine angular
        distribution.

        Default: ESPUTC = 0.
\end{description}

\textbf{Note:} 

All data in a block for ``local reflection data" are, in general, independent
of the type and species of the incident particle.
Some are, however, copied identically NPHOTI + NATMI + NMOLI + NIONI + NPLSI
times onto appropriately dimensioned arrays (with the same names).
This is currently done for the parameters:

\begin{lstlisting}
      ISRS  ISRC
      TRANSP(1, ...)  TRANSP(2, ...)
      RECYCF  RECYCT  RECPRM  EXPPL  EXPEL  EXPIL
      RECYCS  RECYCC  SPTPRM  ESPUTS ESPUTC
\end{lstlisting}

The first index is the species index, the second index is the surface number.

Overwriting local reflection data for some specific incident particle species
can be done via a call to the entry RF0USR of the user supplied reflection
routine REFUSR.
This entry is called in the initialization phase of an EIRENE run, from
subroutine REFLEC.
Likewise, the entry SP0USR of the user supplied sputter routine SPTUSR is
called in the initialization phase from subroutine SPUTER.
For further details on user supplied surface interaction routines see
\cref{sec3.3}.

For example: RECYCT(3,5) is the recycling coefficient for incident species 3
onto surface no.\ 5.

If surface models are defined in input block 6 by the \lstinline|SURFMOD_modname|
label, (rather than individually for each surface in blocks 3a or 3b), then
these variables may be made species dependent by adding an arbitrary number of
lines to a SURFMOD-deck, each of which overwriting the specification for a
particular particle species.
The name of the species (blocks 4 and 5) must be uniquely determining one of
the test particle or bulk species.

In the example of SURFMOD decks given above, e.g., the first such additional
line overwrites the original value of TRANSP(1, \dots) (=0.0) for species D2
with the value 0.5.
I.e., all surfaces, to which the reflection model
modname=\lstinline|CARB_SPT_1153|K is assigned, have a recycling coefficient =1.0
for all incident species, except for the D2 molecules, for which these surfaces
are made transparent with probability 0.5, if these molecules are incident from
the positive side.
For T2 molecules incident from the negative side this transparency is set to
0.25.
The other species dependent modifications in this SURFMOD deck are
self-explaining.

\subsection {effective pumping speed}\label{sec2.6.1}

The flag RECYCT described in the previous section is also used to specify
surface pumping in the following way:

Let $A$ be the surface area [\si{\square\centi\metre}] of a surface (additional
surface or non-default standard surface) as seen by test particles, to which a
given pumping speed $S$ is to be assigned.
Then the pumping speed $S$ [\si{\litre\per\second}] for particles with
temperature $T$ [\si{\kelvin}] and mass $m$ [\si{\atomicmassunit}] is related
to the sticking fraction $p_a = 1 -RECYCT$ by
\begin{equation}
  \label{pumping-speed} S = A \cdot (1 - RECYCT) \cdot 3.638 \cdot \sqrt{T/m} =
  \tilde A \cdot 3.638 \cdot \sqrt{T/m}
\end{equation}
If the surface area is
given in [\si{\square\metre}], $S$ in [\si{\cubic\metre\per\second}], then the
numerical factor becomes \num{36.38}.

Note: the first two factors define an \emph{effective exposed area} $\tilde A$,
across which particles with a thermal velocity would be removed.
The remaining factors are the (thermal) effective velocity of particles lost
across this area: $S = \tilde A \cdot  v_{th}$ [volume /time]

The pump-throughput (pumped flux), (mass flow rate) $Q$ is the product of
pressure $P$  in front of surface $A$ and pumping speed $S$:

$Q_{pump} = P \cdot S$, 

and, for a given type of gas ($m$) and gas temperature $T$  the pump-throughput
is given e.g.\ in flux units \si{\per\second}, (or: \si{\ampere}), or as mass
flow with dimension pressure times volume per time.

The temperature $T$ and type of gas (given by $m$) are part of the
specification of $S$, e.g.\ to be found from the technical specification of the
pump.
$S$ may also depend on the pressure for some pumps: $S = S(P,T,m)$.
However in the molecular flow regime $S$ may often assumed to be a constant
(independent of pressure) for a given temperature and type of gas.

But normally even in this case	the specified pumping speeds $S$ for vacuum
pumps do not follow the $\sqrt{T/m}$ dependence in equation
\cref{pumping-speed} for different types of gas.
E.g.\ pumping speeds for  H$_2$ vs.\ He of ``real pumps", for same
temperatures, are usually not simply related by a factor $1/\sqrt{2}$, and
those for D$_2$ and He may not be identical either.
In a simulation of a gas mixture then a species dependent flag RECYCT (to
produce a species dependent effective area $\tilde{A}$(ISPZ) has to be
specified to reproduce a given ratio of pumping efficiencies for different
types of gas, for a particular type of pump (see the ``Note" at the end of
previous section: RECYCT(ISPZ) option) to reproduce the proper species
dependence (and pressure dependence, if any) of the pumping system.

Probably sometimes, when $S=S(P)$, even RECYCT has to be made dependent on the
pressure $P$ found in front of the surface $A$ (then e.g.\ by an iterative
procedure)

\subsection{Pressure Feedback Loop}

This type of boundary condition modifies the recycling probability 
{\texttt RECYCT} by means of a proportional-integral loop so that the pressure 
in a given cell matches a given reference value. It is activated by setting 
{\texttt ILREF = 4} in the desired reflection model in block 6a of the Eirene 
input file. Two new values, related with the reference cell and pressure, 
have to be provided. The first one in the line:

\begin{verbatim}
ILREF ILSPT ISRS ISRC REFCELL
\end{verbatim}

in which {\texttt REFCELL} is an integer referencing the cell in the Eirene mesh 
in which the pressure will be computed. This value can be easily obtained 
with {\texttt triang} by means of the {\texttt (I)nquire} function. The second 
value is provided in the line:

\begin{verbatim}
RECYCS RECYCC SPTPRM ESPUTS ESPUTC REFPRESS
\end{verbatim}

in which {\texttt REFPRESS} is the reference pressure in $Pa$.

This loop will not provide values of {\texttt RECYCT} larger than $1$ or below $0$.
The coefficients of the proportional-integral loop are hard-coded.
